% Edition 20: French physical 71 bottom--73 end / English 47 bottom--48 end.
% Author text is translated in order; editorial changes are explicitly footnoted.
\source{法稿\texttt{EGA5\_francais.pdf}物理第71页末至73页末；英译\texttt{EGAVaiello.pdf}第47页下部至48页末}

另一方面，在$\operatorname{Grass}^m$上的$X_{\operatorname{Grass}^m}\setminus V''$中，情形正是“第16段补篇”中所考虑的良好情形。相对于基概形$S$，$V^1\setminus V''$确实光滑（因为它在$X$上光滑）；在所讨论的非空、纯维片上，它的相对维数比$\operatorname{Grass}^m$在$S$上的相对维数小一，正如在上面的公式中取$k=1$所见。因此，所引之处的结果适用。特别地，我们得到：$\phi^m$相对于$\operatorname{Grass}^m$的超奇异点集无非是
\[
V''\cup V^2,
\]
其中$V^2$是
\[
V^1\setminus V''\longrightarrow\operatorname{Grass}^m
\]
的分歧子预概形。因此还可以说，在$V''$以外，超奇异零点由至少两个普通奇异零点的合并而产生。（不过，我们并非一定要这样说。）\cmt{V20-FAM001}{法稿第72页使用$V^1$，英译第47页改排为$V'$；此处据法稿统一回$V^1$，与前段相接。$V''$仍为双撇，$V^2$仍为上标2：后者在本句被定义为分歧子预概形，不能认作按交维数定义的$V^{(2)}$。法稿的环境确为$X_{\operatorname{Grass}^m}$，没有将其改成零点概形$X^{(m)}$。新增“非空、纯维片”限制承接前段已修正的维数公式；空退化层不能赋以该有限维数。这一段依赖所引补篇及前面的光滑性分析，原稿没有在此展开一般$m$的全部论证；“合并”是作者的解释性措辞，并不是一个已给出变形族和重数的独立定理。英译在loc. cit.后加“[Latin]”（拉丁语标记），又以“collapsing together”解释“collapsing”；这两项是英译的词语注释。}

这样，我们基本上得到了与8.7的a）、b）相当的结果。利用$\operatorname{Grass}^m$的切丛的显式描述（类似于$m=1$的情形），应当能够给出与8.7的c）相当的条件。这还意味着：对于$V^1\setminus V''$中一个在$\operatorname{Grass}^m$上无分歧的几何点，只要知道它在$\operatorname{Grass}^m$中的像，就能知道它在$P$中的像，前提是前一个像为$V^1$在$\operatorname{Grass}^m$中的闭像的光滑点（这里假设$S$是一个域的谱）。如有需要，我可以给出一个更精确的陈述。\cmt{V20-FAM002}{法稿第72页的手写改订可读为“cela implique encore que”；因此据法稿恢复“这还意味着”，撤去旧D17-M037依据英译“illegible”而保留的文字缺口。这里的“前提是”对应“pourvu que”，不能弱化为仅作解释的“鉴于”。法稿保留“应当能够”和“更精确的陈述”等措辞，故这仍是拟议的推广；不能把字迹恢复当作8.7(c)缺失论证的恢复。末句的第一人称是原作者，不是本译者的承诺。}

一旦这个问题解决，就有推广8.8、8.9、8.11的那些显然推论。无疑，也可以在任意$m$的情形下陈述第8节的其他命题。如果这需要额外的编写工作，即使收入前面的微分方面的讨论，也可以放弃这些推广。第9节的结果也是如此。\cmt{V20-FAM003}{“一旦这个问题解决”保留原文“Une fois ce point acquis”的条件。这里列出的是原作者的写作计划，不能据此认定8.8的工作假设、8.7(c)的缺续或任意$m$的推广已经获得证明。}

至于第10节，其中所研究的情形可按下述方式推广到任意$m$：固定$P$中一个余维为$m+1$的线性子簇$C$，考虑由$P$中包含$C$、余维为$m$的线性子簇$L$组成的射影空间$Q$。$Q$是$\operatorname{Grass}^m$的闭子预概形；特别地，可以形成我们打算研究的$X_Q^{(m)}$。\cmt{V20-FAM004}{法稿第72页末至73页开头清楚写“sous-pr\'esch\'ema ferm\'e”，英译的“close subscheme”只是英文用词错误。若$P=\mathbf P^N_k$且$1\le m\le N-1$，则$Q(C)\simeq\mathbf P^m_k$：在包含$C$的$(m+1)$维线性形式空间中选择一个$m$维子空间。这一参数维数将在本段的反例与末尾引理中使用。}

第一点无论如何都应当写入正文：至少当$C$“正则地”截$X$时，$X_Q^{(m)}\to X$仍是双有理的；确切地说，在这种情形下，$X_Q^{(m)}$典范地同构于沿$X\times_P C$爆破$X$所得的预概形。这个事实的证明，正是借助10.3得到的10.2的证明。\cmt{V20-FAM005}{法稿第73页完整写出中心$X\times_P C$，据此恢复英译中破损的纤维积排法。此处“正则地”带有原稿引号；与本版10.2、10.3相同，已核实的应用范围是局部诺特情形：在中心附近，定义$C$的$m+1$个线性形式拉回并局部平凡化后构成正则序列；在空中心上采用真空约定。不能删去正则性条件而对任意中心断言上述关联概形就是爆破。一般非诺特正则性的约定仍待统一，原稿未在本段另作证明。}

第二点有一定意义，但我们并非非把它包括进来不可：如果$C$选得“足够一般”，那么$X_Q^{(m)}\to Q$具有某些良好性质。原稿在这里拟议如下推广。设$X$在$S=\Spec(k)$上光滑、固有、几何不可约，维数$n\ge m$。令$T$为使$X_\xi^{(m)}$不是$k(\xi)$上维数为$n-m$的光滑概形的那些$\xi\in Q$组成的集合。\cmt{V20-FAM006}{法稿第73页清楚写$S=\Spec(k)$；英译漏掉$(k)$，故旧D17-M041所留的基域缺口现在可以恢复。关于每个参数处的条件，法稿前两处未添纤维下标，后文则明写$X_\xi^{(m)}$；正文据同段的对象关系补明下标。法稿也把全局集合$T$写成“在$k(\xi)$上几何不可约”，这个移动的剩余域不能作为整个$T$的固定基域；即使订为$k$，该不可约性断言仍有实质反例，见下一注。}

\noindent\textbf{译校订正。}仅凭本段写出的条件，即使加入原作者稍后补充的无分歧和高次Veronese嵌入要求，也不能断言$T$几何不可约。\cmt{V20-FAM007}{原稿紧接上述定义断言：“那么，对足够一般的$C$，$T$在$k(\xi)$上几何不可约，并在$Q$中余维为1。”这句话的不可约性部分不能保留为正文定理。取代数闭的特征零域$k$，$X=\mathbf P^1_k$，将恒等嵌入作二次Veronese变换，得到$[u:v]\mapsto[u^2:uv:v^2]\in\mathbf P^2_k$；令$m=n=1$、$C=[0:1:0]$。经过$C$的线束为$a x_0+c x_2=0$，拉回为$au^2+cv^2=0$，故$T=V(ac)\subset\mathbf P^1_{[a:c]}$由两个不同的点组成，并非几何不可约。对任意位于圆锥曲线外的一般中心，同样有两条不同切线，所以这不是一个特殊中心造成的反例。$X$光滑、固有、几何不可约，给定映射是闭浸入、因而无分歧，且已经作了次数2的Veronese变换。完整核验见报告V20-FAM-E02。原稿的余维和普通奇点方面的其余拟议结论仍逐句保留于下文，但不因删除这一错误结论就一概视为已证。}

原作者拟议的其余结论是：对“足够一般”的$C$，$T$在$Q$中余维为1；令$T^1$为$T$中使$X_\xi^{(m)}$至少有一个“超奇异”点的那些$\xi$组成的集合，则$T^1$在$T$中无处稠密。最后，若$f\colon X\to P$是浸入，那么略微扩大$T^1$后，对每个$\xi\in T\setminus T^1$，$X_\xi^{(m)}$中恰有一个非光滑点，而且该点在$k(\xi)$上有理。\cmt{V20-FAM008}{法稿手添的例外集记号为$T^1$，英译改排$T'$；此处据法稿恢复。法稿的态射记号为小写$f$，英译写成大写$F$，不能把这一误排归为法稿原错。上一句以“原作者拟议的其余结论是”明示本版编排：保留原文余维、稀疏性及唯一奇点的全部叙述，但不将它们包装成本轮已完成一般证明的新定理。}

我忘了在陈述中说明：假设$X\to P$无分歧，而且必须预先把$f$替换为$\phi_d f$，其中$d\ge2$（$\phi_d$在9.1中定义）。\cmt{V20-FAM009}{原稿在Veronese指标处也使用$n$，与本段已经表示$\dim X$的$n$重名。正文仅将这个指标改名为$d$，原式为$\phi_n f$、$n\ge2$；不把两个用途强行认作同一个数。这一换名避免在曲线情形中误读出“$\dim X\ge2$”的额外条件，原稿没有该限制。原作者后补的两个假设是整段条件的一部分，不能遗漏；它们仍不能排除上一注的二次Veronese圆锥曲线反例。}

证明这一陈述最自然的办法，似乎是利用$\operatorname{Grass}^m$中对应于“奇异”截面$X_\xi^{(m)}$的子概形$Z$（在8.8中记为$T$）。按照拟议中的前述推广，在所说的条件下，它几何不可约，余维为1，而对应于“超奇异”截面的子概形$Z^1$在其中无处稠密。\cmt{V20-FAM010}{法稿仍是“le plus naturel pour prouver cet \'enonc\'e semble \^etre…”的证明设想，并写“on sait, sous les conditions dites”。正文补明“按照拟议中的前述推广”，保留其依赖关系；没有在本轮把任意$m$的这些判别概形性质另作完整证明。即使已知$Z$几何不可约，也不能直接推出$Z\cap Q(C)$几何不可约；这正是下一句拟议引理的逻辑缺口。}

于是，还须证明一个大致如下的引理：设$Z$是$\operatorname{Grass}^m$中余维为$q$的闭子集，并如上由$C$定义$Q(C)$。那么，对于每个“足够一般”的$C$，交集
\[
Q(C)\cap Z
\]
在$Q(C)$中的余维至少为$q$。\cmt{V20-FAM011}{法稿第73页末的交集明确写$Q(C)\cap Z$；英译第48页第一处误排为$Q(Z)\cap Z$。旧D17-M044的数学订正方向正确，本轮进一步把错误归属限定为英文转写。在域上、通常的非空中心范围$1\le m\le\dim P-1$内，这一余维界可由二步旗概形和一般纤维维数证明；报告V20-FAM-E03给出独立补证，空交约定余维为无穷。这个补证不属于原稿。有限域上还须区分非空开参数集与其中是否存在基域有理点，不能由“足够一般”自动断言后者。}

\noindent\textbf{译校订正。}这一余维界本身不包含交集的几何不可约性；几何不可约的$Z$与一般$Q(C)$相交，仍可能可约。\cmt{V20-FAM012}{法稿末行的手写补句可辨为：“此外，若$Z$几何不可约，$Q(C)\cap Z$也同样如此，如果$Z$‘足够一般’。”其中“itou”意为“同样”，最后一个字母仍是$Z$；不擅改成$C$。英译在这里留下“illegible, ask AG”及“itou???”，旧D17-M045因此把谓语记为残缺；法稿现可恢复拟议的不可约性结论，然而这个结论不能照录为真命题。上面的圆锥曲线反例中，$Z$是对偶平面中光滑、几何不可约的对偶圆锥曲线，$Q(C)$是一般直线，交集为两个不同的点；若“$Z$足够一般”指固定次数的光滑圆锥曲线，一般的这样的$Z$仍给出反例。原稿没有指定能保证该断言的更窄参数空间或额外条件，故保留其原说法于本注，不补造正确的一般不可约性定理。第12节止于这一手写补句；下一原页开始第13节。}
